\documentclass[11pt]{article}
%% Packages for bibliography
% Use either natbib or biblatex
%\usepackage[round]{natbib}  

\usepackage[backend=biber,
            bibstyle=authoryear,
            citestyle=authoryear,
            maxbibnames=9,
            maxcitenames=2,
            dashed=false,
            firstinits=true,]
            {biblatex}
 
\usepackage{bibentry}

% See other styles in sharelatex
% https://www.sharelatex.com/learn/Biblatex_bibliography_styles


%% Hyperref authoryear style
% https://tex.stackexchange.com/questions/1687/hyperlink-name-with-biblatex-authoryear
\DeclareCiteCommand{\cite}
  {\usebibmacro{prenote}}
  {\usebibmacro{citeindex}%
   \printtext[bibhyperref]{\usebibmacro{cite}}}
  {\multicitedelim}
  {\usebibmacro{postnote}}

\DeclareCiteCommand*{\cite}
  {\usebibmacro{prenote}}
  {\usebibmacro{citeindex}%
   \printtext[bibhyperref]{\usebibmacro{citeyear}}}
  {\multicitedelim}
  {\usebibmacro{postnote}}

\DeclareCiteCommand{\parencite}[\mkbibparens]
  {\usebibmacro{prenote}}
  {\usebibmacro{citeindex}%
    \printtext[bibhyperref]{\usebibmacro{cite}}}
  {\multicitedelim}
  {\usebibmacro{postnote}}

\DeclareCiteCommand*{\parencite}[\mkbibparens]
  {\usebibmacro{prenote}}
  {\usebibmacro{citeindex}%
    \printtext[bibhyperref]{\usebibmacro{citeyear}}}
  {\multicitedelim}
  {\usebibmacro{postnote}}

\DeclareCiteCommand{\footcite}[\mkbibfootnote]
  {\usebibmacro{prenote}}
  {\usebibmacro{citeindex}%
  \printtext[bibhyperref]{ \usebibmacro{cite}}}
  {\multicitedelim}
  {\usebibmacro{postnote}}

\DeclareCiteCommand{\footcitetext}[\mkbibfootnotetext]
  {\usebibmacro{prenote}}
  {\usebibmacro{citeindex}%
   \printtext[bibhyperref]{\usebibmacro{cite}}}
  {\multicitedelim}
  {\usebibmacro{postnote}}

\DeclareCiteCommand{\textcite}
  {\boolfalse{cbx:parens}}
  {\usebibmacro{citeindex}%
   \printtext[bibhyperref]{\usebibmacro{textcite}}}
  {\ifbool{cbx:parens}
     {\bibcloseparen\global\boolfalse{cbx:parens}}
     {}%
   \multicitedelim}
  {\usebibmacro{textcite:postnote}}
  
  
  
  
%% Packages for figures
\usepackage{graphicx}			%% Path to figures and more I guess
\usepackage{epstopdf}			%% Needed for .eps files
\usepackage{subfig}		    	%% Array of figures (incompatible with subcaption)
\usepackage{float} 				%% Force floats position


%% Packages for tables
\usepackage{tabularx}			%% Advanced tables
\usepackage{booktabs}			%% http://www.tablesgenerator.com/
\usepackage{multirow}			%% Combine cells of tables
\usepackage[flushleft]{threeparttable} % Add notes to a table
\usepackage{adjustbox} 			%% Give maximum height-width
\usepackage{rotating}			%% Sideways table
%\usepackage{floatrow}			%% Table and figure side by side


%% Packages for captions
\usepackage[				    	%% Caption style
font=small,labelfont=bf]
{caption}
%\usepackage{subcaption}			%% Make captions for subfigures
\captionsetup{parskip=5pt}	    	%% Set the caption paragraph spacing


%% Packages for math
\usepackage{amssymb}
\usepackage{mathrsfs}
\usepackage{amsthm}
\usepackage[fleqn]{amsmath}		%% Subequations and more I guess
\usepackage{nicefrac}			%% Cool fractions
\usepackage[fleqn]{cases}		%% Piecewise function
\usepackage{bm}					%% Bold greek symbols in math mode
\usepackage{dsfont}             %% Caligraphic math months
\usepackage{commath}            %% \norm


%% Packages for links
\usepackage{color}				%% Color for the text
\usepackage{xcolor}				%% Advanced color for the text
\usepackage[					%% Clickable links
hypertexnames=false,
bookmarksopen=true]{hyperref}	
\hypersetup{
    colorlinks,
    linkcolor={blue!00!black},
    citecolor={blue!30!black},
    urlcolor={blue!30!black}}
\usepackage{url}				%% Hyperlinks to webs
% The option hypertexnames=false is set to prevent problems when the chapter numbering is modified (appendices)


%% Dots for the table of contents
\usepackage{tocloft}
\renewcommand{\cftpartleader}{\cftdotfill{\cftdotsep}} % for parts
\renewcommand{\cftchapleader}{\cftdotfill{\cftdotsep}} % for chapters
\renewcommand{\cftsecleader}{\cftdotfill{\cftdotsep}} % for sections, if you really want! (It is default in report and book class (So you may not need it).
% https://tex.stackexchange.com/questions/53898/how-to-get-lines-with-dots-in-the-table-of-contents-for-sections
% ----------------------------------------------------------------


%% Package for nomenclature
\usepackage[intoc]{nomencl}
\makenomenclature
% Create nomenclature groups
\usepackage{etoolbox}
\renewcommand\nomgroup[1]{%
  \item[\large\bfseries
  \ifstrequal{#1}{A}{Symbols}{%
  \ifstrequal{#1}{B}{Greek Letters}{%
  \ifstrequal{#1}{C}{Acronyms}{%
  \ifstrequal{#1}{D}{Subscripts}{}}}}%
]\vspace{12pt}}
% Add the units to nomenclature
\newcommand{\nomunit}[1]{%
\renewcommand{\nomentryend}{\hspace*{\fill}#1}}
% The difficult thing is to prepare the compilation order
% Check the documentation in the future


%% Package for glossary
\usepackage[toc,nonumberlist]{glossaries}
%\setglossarystyle{long}
\setglossarystyle{altlist}
\renewcommand{\glsnamefont}[1]{\textbf{#1}}

\setlength\LTleft{0pt}
\setlength\LTright{0pt}


\makeglossaries
% The difficult thing is to prepare the compilation order
% It requires a Perl distribution to work


%% Packages to create lists (see my equations)
\usepackage{enumitem}          %% List of definitions

%% Spacing before a list and among the elements of the list
\setlist[itemize]{itemsep=1pt, topsep=1pt}
\setlist[enumerate]{itemsep=1pt, topsep=1pt}



%% Packages to create appendices
\usepackage{appendix}			%% Create appendix
\usepackage{ifthen}     		%% Used to remove appendix tables and figures from the list of tables and figures



%% Type of letter
\usepackage{times}


%% Set the indentation space
\setlength{\parindent}{5pt}

%% Set the spacing between paragraphs
\setlength{\parskip}{5pt}

%\usepackage{parskip}   			%% No idea what this does


%% Set the spacing before and after equations
\expandafter\def\expandafter\normalsize\expandafter{%
    \normalsize
    \setlength\abovedisplayskip{6pt}
    \setlength\belowdisplayskip{6pt}
    \setlength\abovedisplayshortskip{6pt}
    \setlength\belowdisplayshortskip{6pt}
}


%% Set the spacing before and after title sections
\usepackage{titlesec}
\titlespacing\section{0pt}{18pt plus 0pt minus 0pt}{2pt plus 0pt minus 0pt}
\titlespacing\subsection{0pt}{12pt plus 0pt minus 0pt}{2pt plus 0pt minus 0pt}
\titlespacing\subsubsection{0pt}{8pt plus 0pt minus 0pt}{0pt plus 0pt minus 0pt}




% %% Set the spacing for the table of contents
% \usepackage{setspace}

% %% Remove chapter spacing from lof and lot
% \usepackage{etoolbox}% http://ctan.org/pkg/etoolbox
% \makeatletter
% % \patchcmd{<cmd>}{<search>}{<replace>}{<succes>}{<failure>}
% \patchcmd{\@chapter}{\addtocontents{lof}{\protect\addvspace{10\p@}}}{}{}{}
% \patchcmd{\@chapter}{\addtocontents{lot}{\protect\addvspace{10\p@}}}{}{}{}
% \makeatother

% %% Packages specific from this template
% \usepackage[Lenny]{fncychap}	%% Makes cool chapter headings

% \usepackage{fancyhdr}
% \newcommand{\HRule}{\rule{\linewidth}{0.5mm}}
% \renewcommand*\contentsname{Table of Contents}
% \pagestyle{fancy}
% \fancyhf{}
% \renewcommand{\chaptermark}[1]{\markboth{\chaptername\ \thechapter.\ #1}{}}
% \renewcommand{\sectionmark}[1]{\markright{\thesection\ #1}}
% \renewcommand{\headrulewidth}{0.1ex}
% \renewcommand{\footrulewidth}{0.1ex}
% \fancypagestyle{plain}{\fancyhf{}\fancyfoot[LE,RO]{\thepage}\renewcommand{\headrulewidth}{0ex}}

% %% Reduce the spacing before the chapter
% \makeatletter
% \patchcmd{\@makechapterhead}{\vspace*{50\p@}}{\vspace*{-20\p@}}{}{}
% \patchcmd{\@makeschapterhead}{\vspace*{50\p@}}{\vspace*{-20\p@}}{}{}
% \patchcmd{\DOTI}{\vskip 80\p@}{\vskip 40\p@}{}{}
% \patchcmd{\DOTIS}{\vskip 40\p@}{\vskip 0\p@}{}{}
% \makeatother
% %https://tex.stackexchange.com/questions/111643/decrease-space-before-and-after-chapter-in-fncychaps

% %% Add the back page
% \usepackage{afterpage}

% \newcommand\blankpage{%
%     \null
%     \thispagestyle{empty}%
%     \addtocounter{page}{-1}%
%     \newpage}
    
%%% Prepare the numbering for the first sections
% \pagenumbering{arabic} 				
% \setcounter{page}{1}
% \pagestyle{fancy}
% \fancyhf{}
% \renewcommand{\chaptermark}[1]{\markboth{\chaptername\ \thechapter.\ #1}{}}
% \renewcommand{\sectionmark}[1]{\markright{\thesection\ #1}}
% \renewcommand{\headrulewidth}{0.1ex}
% \renewcommand{\footrulewidth}{0.1ex}
% \fancyfoot[LE,RO]{\thepage}
% \fancypagestyle{plain}{\fancyhf{}\fancyfoot[LE,RO]{\thepage}\renewcommand{\headrulewidth}{0ex}}

% % Protect the table of contents stretching
% \addtocontents{toc}{\protect\setstretch{0.8}}




% Avoid footnote page breaks
\interfootnotelinepenalty=10000




\usepackage[a4paper,width=165mm,top=20mm,bottom=20mm,bindingoffset=0mm]{geometry}

\renewcommand*\rmdefault{cmr}  % Computer Modern Roman
% \renewcommand*\rmdefault{lmr}  % Latin Modern Roman
% \renewcommand*\rmdefault{bch}  % Charter
% \renewcommand*\rmdefault{ppl}  % Palatino
% \renewcommand*\rmdefault{ptm}  % Times

% Ref sections
\newcommand{\RefSection}[1]{\hyperref[#1]{Section~\ref{#1}}}
\newcommand{\RefSections}[2]{\hyperref[#1]{Sections~\ref{#1}~\textcolor{black}{and}~\ref{#2}}}
\newcommand{\RefAppendix}[1]{\hyperref[#1]{Appendix~\ref{#1}}}
% Ref figures
\newcommand{\RefFig}[1]{\hyperref[#1]{Figure~\ref{#1}}}
\newcommand{\RefFigs}[2]{\hyperref[#1]{Figures~\ref{#1}~\textcolor{black}{and}~\ref{#2}}}
\newcommand{\RefFigss}[2]{\hyperref[#1]{Figures~\ref{#1}~\textcolor{black}{to}~\ref{#2}}}
% Ref tables
\newcommand{\RefTable}[1]{\hyperref[#1]{Table~\ref{#1}}}
\newcommand{\RefTables}[2]{\hyperref[#1]{Tables~\ref{#1}~\textcolor{black}{and}~\ref{#2}}}
\newcommand{\RefTabless}[2]{\hyperref[#1]{Tables~\ref{#1}~\textcolor{black}{to}~\ref{#2}}}
% Ref equations
\newcommand{\RefEq}[1]{\hyperref[#1]{Eq.~(\ref{#1})}}
\newcommand{\RefEqs}[2]{\hyperref[#1]{Eqs.~(\ref{#1})~\textcolor{black}{and}~(\ref{#2})}}
\newcommand{\RefEqss}[2]{\hyperref[#1]{Eqs.~(\ref{#1})~\textcolor{black}{to}~(\ref{#2})}}


\newcommand{\Bezier}[0]{B\'{e}zier }

% Load bibliographic file
\addbibresource{references.bib}


\begin{document}
\thispagestyle{empty}

\begin{center}
{\LARGE \bf A briefing on NURBS curves and surfaces}\\
\vspace{2ex}
{\large Roberto Agromayor, \today }\\
\end{center}

This note introduces NURBS curves and surfaces and summarizes their most important mathematical properties. Most of the material used to prepare this note is available in the NURBS book \parencite[Chapter 4]{NURBS_book}.

\section{NURBS curves}

%% --------------------------------------------------------------------------------------------- %%
%% --------------------------------------------------------------------------------------------- %%
%% --------------------------------------------------------------------------------------------- %%
\subsection{Definition}

A Non-Uniform Rational Basis Spline (NURBS) curve is a parametric curve defined by:
\begin{align}
    \label{eq:def_nurbs_curve}
    \mathbf{C}(u)=\sum_{i=0}^{n} R_{i,p}(u) \, \mathbf{P}_{i} \quad \quad 0 \leq u \leq 1
\end{align}

where $p$ is the order of the curve, the coefficients $\mathbf{P}_{i}$ are called control points, and $R_{i,p}$ are the rational basis functions given by:
\begin{align}
    R_{i,p}(u) = \frac{N_{i,p}(u) \, w_{i}}{\sum\limits_{k=0}^{n} N_{k,p}(u) \, w_{k}}
\end{align}

in which $w_{i}$ are the weights of the control points and $N_{i,p}$ are B-Spline basis functions defined on the non-decreasing knot vector $U$:
\begin{align}
    U = [u_0,\, \dots,\, u_{r}] \in \mathds{R}^{r+1} \quad \text{with} \quad r = n + p + 1
\end{align}

The B-Spline basis functions are given by the recursive relation:
\begin{align}
    N_{i,0}(u) =  \begin{cases}
    1& \mathrm{if} \quad u_{i}\leq u<u_{i+1}\\
    0& \mathrm{otherwise}
    \end{cases}
\end{align}
\begin{align}
 N_{i,p}(u) = \frac {u-u_{i}}{u_{i+p}-u_{i}} N_{i,p-1}(u) + \frac {u_{i+p+1}-u}{u_{i+p+1}-u_{i+1}} N_{i+1,p-1}(u)
\end{align}

This recursive relation can yield the $0/0$ quotient which is defined to be zero by convention.


%% --------------------------------------------------------------------------------------------- %%
%% --------------------------------------------------------------------------------------------- %%
%% --------------------------------------------------------------------------------------------- %%
\subsection{Mathematical properties of rational basis functions}
Here is a list of some important properties of the rational basis functions:

\begin{enumerate}

    \item The relation $r = n + p + 1$ holds, where $n+1$ is the number of basis functions and $r+1$ is the number of elements of the knot vector $U$.

    \item The numerator and denominator of $R_{i,p}(u)$ are, at most, polynomials of degree $p$.

    \item Local support:
    \begin{enumerate}
        \item $R_{i,p}(u)=0$ if $u$ is outside the interval $[u_{i}, u_{i+p+1})$. 
        \item In any given knot interval, $[u_{i_{0}}, u_{i_{0}+1})$, at most $(p+1)$ basis functions are nonzero, namely the $R_{i,p}(u)$ with $i_{0}-p \leq i \leq i_{0}$.
    \end{enumerate}

    \item Non-negativity:
        \begin{equation*}
            R_{i,p}(u) \geq 0 \quad \text{for all $i$, $p$, and $u \in [0, 1]$}
        \end{equation*}
        
    \item Partition of unity:
        \begin{equation*}
            \sum_{i=0}^n R_{i,n}(u) = 1 \text{ for all } u \in [0, 1]
        \end{equation*}
    
    In addition, for an arbitrary knot span, $[u_{i_{0}}, u_{i_{0}+1})$ we have that:
        \begin{equation*}
            \sum_{i=i_{0}-p}^{i_{0}} R_{i,n}(u) = 1 \text{ for all } u \in [u_{i_{0}}, u_{i_{0}+1})
        \end{equation*}
    This means that the sum of the non-zero basis-functions of any knot span is unity.

    \item Continuity and differentiability:
        \begin{enumerate}
            \item The basis functions are infinitely differentiable in the interior of the knot intervals.
            \item The basis functions are $p-k$ continuously differentiable at a knot with multiplicity $k$.
        \end{enumerate}
        
    \item Extrema:
    
    $R_{i,p}(u)$ attains exactly one maximum value in the interval $u \in [0, 1]$.
    
    \item There is no compact and useful expression for the first derivative of the basis functions.
    
    \item There is no compact and useful expression for the $k$-th order derivative of the basis functions.
        
    \item When the first and last knots have multiplicity $p+1$ the knot vector is given by:
        \begin{align*}
            & U = [\underbrace{u_0,\, u_1,\, \dots,\, u_{p}}_{p+1},\, \underbrace{u_{p+1},\, \dots,\, u_{n}}_{n-p},\, \underbrace{u_{n+1},\, \dots,\, u_{n+p},\, u_{n+p+1}}_{p+1}]
        \end{align*}
        where:
        \begin{align*}
            & u_0 = \dots = u_p = 0 \\
            & u_{n+1} = \dots = u_{n+p+1} = 1
        \end{align*}
    and is called a \textit{clamped knot vector}. Basis functions of clamped knot vectors satisfy two additional properties:
    \begin{enumerate}
            \item $R_{0,p}(u=0) =1$ and $R_{i,p}(u=0)=0$ for $i\neq0$
            \item $R_{n,p}(u=1) =1$ and $R_{i,p}(u=1)=0$ for $i\neq n$
    \end{enumerate}
    
    \item If all the control point weights are equal, then rational basis functions reduce to B-Spline basis functions, that is, $R_{i,p}(u)=N_{i,p}(u)$
    
    \item If all the control point weights are equal, the knot vector is clamped and $p=n$, then rational basis function reduce to Bernstein polynomials, that is, $R_{i,p}(u)=B_{i,n}(u)$
    
\end{enumerate}




%% --------------------------------------------------------------------------------------------- %%
%% --------------------------------------------------------------------------------------------- %%
%% --------------------------------------------------------------------------------------------- %%
\subsection{Mathematical properties of NURBS curves}
Here is a list of some important properties of NURBS curves:

\begin{enumerate}

    \item $\mathbf{C}(u)$ is a piecewise curve and its components are ratios of polynomials of, at most, degree $p$.
    
    \item The order $p$, number of control points $n+1$ and number of knots $r+1$ are related according to $r = n + p + 1$.
    
    \item NURBS curves contain B-Spline and rational/non-rational \Bezier curves as special cases:
    \begin{enumerate}
        \item If all the control point weights are equal then the NURBS curve is reduced to a B-Spline curve.
        \item If $n = p$ and the knot vector is clamped then the NURBS curve is reduced to a rational \Bezier curve.
        \item If all the control point weights are equal, $p = n$, and the knot vector is clamped then the NURBS curve is reduced to a polynomial \Bezier curve.
    \end{enumerate}

    \item Affine invariance:
    
    NURBS curves are invariant under affine transformations such as rotations, displacements and scalings. This means that one can apply an affine transformation to the NURBS curve by applying it to its set of control points.
    
    Given an affine transformation $\phi(\mathbf{v}) = \mathrm{A} \mathbf{v} + \mathbf{b}$ we have that:
    \begin{align*}
        \phi \big(\mathbf{C}(u)\big) &= \phi \big(\sum_{i=0}^{n} R_{i,p}(u)\, \mathbf{P}_{i}\big) = \\
        &= \mathrm{A} \sum_{i=0}^{n} R_{i,p}(u)\, \mathbf{P}_{i} + \mathbf{b} = \\
        &= \mathrm{A} \sum_{i=0}^{n} R_{i,p}(u)\, \mathbf{P}_{i} + \mathbf{b} \, \sum_{i=0}^{n} R_{i,p}(u) \qquad \text{(by partition of unity)}\\
        &= \sum_{i=0}^{n} R_{i,p}(u)\,( \mathrm{A} \mathbf{P}_{i} + \mathbf{b}) \\
        &= \sum_{i=0}^{n} R_{i,p}(u)\,\phi \big(\mathbf{P}_{i} \big)
    \end{align*}
    
    \item Convex hull property:
    
    All the points of a NURBS curve are contained in the \textit{convex hull} of its control points.
    The convex hull of a set of points $P = \{\mathbf{P}_{0},\, \dots,\, \mathbf{P}_{N}\}$ is denoted by $\mathcal{CH}(P)$ and it is the set of all the convex combinations of points:
    \begin{equation*}
     \mathcal{CH}(P) = \Big\{ \mathbf{C} =  \sum_{k=0}^{N} a_{k} \, \mathbf{P}_{k} \text{ such that } \sum_{k=0}^{N} a_{k} = 1 \text{ and } a_{k}\geq 0 \text{ for } k = 0,\, 1,\, \dots,\, N\Big\}
    \end{equation*}
    The convex hull property follows from the non-negativity and the partition-of-unity properties of the basis functions.
    % The convex hull of a set of points is the minimal convex polytope that contains all the points
    
    \item Strong convex hull property:
    
    If $u \in [u_{i_{0}},\, u_{i_{0}+1})$, then $\mathbf{C}(u)$ is contained in the convex hull of the control points $\boldsymbol{P}_{i}$ with $i_{0}-p \leq i \leq i_{0}$.

    \item Local modification scheme:
    
    Modifying the control point $\mathbf{P}_{i}$ or weight $w_{i}$ affects $\mathbf{C}(u)$ only in the interval $[u_{i}, u_{i+p+1})$. This property follows from the fact that $R_{i,p}(u)=0$ if $u$ is outside the interval $[u_{i}, u_{i+p+1})$ and it implies that the shape of a NURBS can be modified locally without changing its shape globally.
    
    \item The polygon formed by the set of control points is known as \textit{control polygon}. The control polygon represents a piecewise linear approximation to the NURBS curve.

    \item Variation diminishing property:
    
    No straight line (or plane in three dimensions) intersects the B-Spline curve more times than it intersects its control polygon. An intuitive explanation of this property is that the B-Spline curve does not wiggle more than its control polygon.
    
    \item Continuity and differentiability:
    \begin{enumerate}
        \item NURBS curves are infinitely differentiable in the interior of the knot intervals.
        \item NURBS curves are \textit{at least} $p-k$ continuously differentiable at a knot with multiplicity $k$.
    \end{enumerate}
    
    
    \item Homogeneous coordinates:
    
    $N$-dimensional rational functions with a common denominator can be represented as polynomial functions in $(N+1)$-dimensional space using \textit{homogeneous coordinates}.
    
    Consider a three-dimensional point $\mathbf{P}=(x,\, y,\, z)$. $\mathbf{P}$ can be written in four-dimensional space as $\mathbf{P}^w=(wx,\, wy,\, wz,\, w) = (X,\, Y,\, Z,\, W)$, where $w \neq 0$.
    The point $\mathbf{P}$ can be obtained from $\mathbf{P}^w$ using the mapping $\mathcal{H}$ given by:
    \begin{align*}
        \mathbf{P} = \mathcal{H}\{\mathbf{P}^w\} =  \mathcal{H}\{(X,\, Y,\, Z,\, W)\} = \left( \frac{X}{W},\, \frac{Y}{W},\, \frac{Z}{W}\right)
    \end{align*}
    
    This theory can be used to represent a $N$-dimensional NURBS curve as a $(N+1)$-dimensional B-Spline curve and then retrieve the coordinates of the NURBS curve using the $\mathcal{H}$ mapping.
    
    Indeed, consider a NURBS curve $\mathbf{C}(u)$ with control points $\mathbf{P}_{i}$ and weights $w_{i}$. It is possible to construct a set of weighted control points $\mathbf{P}^{w}_{i} = (w_{i}x_{i},\, w_{i}y_{i},\, w_{i}z_{i},\, w_{i})$ and define the corresponding non-rational B-Spline curve in four-dimensional space as:
    \begin{align*}
        \mathbf{C}^{w}(u)=\sum_{i=0}^{n} N_{i,p}(u)\, \mathbf{P}^{w}_{i} \quad \quad 0 \leq u \leq 1
    \end{align*}
    
    In order to compute the coordinates of the original NURBS curve it is possible to use standard B-Spline algorithms to evaluate the coordinates of $\mathbf{C}^{w}(u)$ in homogeneous space and then apply the $\mathcal{H}$ mapping:
    \begin{align*}
        \mathbf{C}(u) &= \mathcal{H}\{\mathbf{C}^w\} \\
        &=  \mathcal{H} \bigg\{ \sum_{i=0}^{n} N_{i,p}(u)\, \mathbf{P}^{w}_{i} \bigg\} \\
        &= \frac{\sum\limits_{i=0}^{n} N_{i,p}(u)\, w_{i} \,\mathbf{P}_{i}}{\sum\limits_{i=0}^{n} N_{i,p}(u)\, w_{i}} \\
        &= \sum_{i=0}^{n} R_{i,p}(u) \, \mathbf{P}_{i}
    \end{align*}
    
    This property is important to evaluate the derivatives of NURBS curves.
    
    \item First and higher order derivatives:
    
    The derivatives of a NURBS curve $\mathbf{C}(u)$ can be expressed in terms of the derivatives of the corresponding B-Spline in homogeneous space $\mathbf{C}^{w}(u)$. Indeed, let:
    \begin{align*}
        \mathbf{C}(u) = \frac{\sum\limits_{i=0}^{n} N_{i,p}(u)\, w_{i} \,\mathbf{P}_{i}}{\sum\limits_{i=0}^{n} N_{i,p}(u)\, w_{i}} = \frac{\mathbf{A}(u)}{w(u)}
    \end{align*}
    
    where $\mathbf{A}(u)$ is the vector function whose coordinates are the three first elements of $\mathbf{C}^{w}(u)$ and $w(u)$ is the scalar function with the fourth element of $\mathbf{C}^{w}(u)$.
    
    To compute the first derivatives, first rearrange the previous expression to avoid the denominator, then differentiate both sides of the equation and, finally, solve for $\mathbf{C}'(u)$:
    \begin{align*}
        & w \, \mathbf{C} = \mathbf{A} \\
        & \big[ w \, \mathbf{C} \big] ' = \mathbf{A}' \\
        & w' \, \mathbf{C} + w \, \mathbf{C}' = \mathbf{A}' \\
        & \mathbf{C}' = \frac{1}{w} \left(\mathbf{A}' - w' \, \mathbf{C} \right)
    \end{align*}
    
    To compute the $k$-th order derivatives, differentiate $k$-times using Leibniz' rule for product differentiation and then solve for $\mathbf{C}^{(k)}(u)$:
    \begin{align*}
        & w \, \mathbf{C} = \mathbf{A} \\
        & \big[ w \, \mathbf{C} \big]^{(k)} = \mathbf{A}^{(k)} \\
        & \sum_{i=0}^{k} \binom{k}{i} w^{(i)} \, \mathbf{C}^{(k-i)} = \mathbf{A}^{(k)} \\
        & \mathbf{C}^{(k)} =  \frac{1}{w} \left(\mathbf{A}^{(k)} - \sum\limits_{i=1}^{k} \binom{k}{i} w^{(i)} \, \mathbf{C}^{(k-i)} \right)
    \end{align*}
    
    The derivatives of $\mathbf{A}(u)$ and $w(u)$ can be easily computed by differentiating $\mathbf{C}^{w}(u)$:
        \begin{align*}
             \left[\mathbf{C}^{w}\right]^{(k)}(u)  = \sum_{i=0}^{n}N^{(k)}_{i,p}(u)\, \mathbf{P}^{w}_{i}
        \end{align*}
    

    \item End point interpolation:
    
    The start and end points of a clamped NURBS curve coincide with the first and last control points, respectively.
    \begin{align*}
        \mathbf{C}(u=0) &= \mathbf{P}_{0} \\
        \mathbf{C}(u=1) &= \mathbf{P}_{n}
    \end{align*}
    
    
    \item End point tangency:
    
    A clamped NURBS curve is tangent to the control polygon at the endpoints.
    \begin{align*}
        \frac{\mathrm{d}\mathbf{C}}{\mathrm{d}u}\Bigr|_{u=0} &= \left(\frac{p}{u_{p+1}}\right) \left(\frac{w_{1}}{w_{0}}\right) (\mathbf{P}_{1} - \mathbf{P}_{0}) \\
        \frac{\mathrm{d}\mathbf{C}}{\mathrm{d}u}\Bigr|_{u=1} &= \left(\frac{p}{1-u_{n}}\right) \left(\frac{w_{n-1}}{w_{n}}\right) (\mathbf{P}_{n} - \mathbf{P}_{n-1})
    \end{align*}
    
    \item End point curvature:
    
    The second derivative of NURBS curve at its endpoints in given by: \\
    \resizebox{\linewidth}{!}{
    \begin{minipage}{0.99\linewidth}
    \begin{align*}
        \frac{\mathrm{d}^2\mathbf{C}}{\mathrm{d}u^2}\Bigr|_{u=0} =& \frac{p (p-1)}{u_{p+1}}  \left[ \left(\frac{1}{u_{p+2}} \right) \left(\frac{w_{2}}{w_{0}} \right) \left(\mathbf{P}_{2} - \mathbf{P}_{0} \right) - \left( \frac{1}{u_{p+1}}  + \frac{1}{u_{p+2}} \right) \left(\frac{w_{1}}{w_{0}} \right) \left( \mathbf{P}_{1} - \mathbf{P}_{0} \right) \right] + \frac{2p^{2}}{u^{2}_{p+1}} \left(\frac{w_{1}}{w_{0}} \right) \left(1 - \frac{w_{1}}{w_{0}} \right) \left( \mathbf{P}_{1} - \mathbf{P}_{0} \right) \\
        \frac{\mathrm{d}^2\mathbf{C}}{\mathrm{d}u^2}\Bigr|_{u=1} =& \frac{p (p-1)}{1-u_{n}}  \left[ \left(\frac{1}{1- u_{n-1}} \right) \left(\frac{w_{n-2}}{w_{n}} \right) \left(\mathbf{P}_{n-2} - \mathbf{P}_{n} \right) - \left( \frac{1}{1-u_{n}}  + \frac{1}{1-u_{n-1}} \right) \left(\frac{w_{n-1}}{w_{n}} \right) \left( \mathbf{P}_{n-1} - \mathbf{P}_{n} \right) \right] + \frac{2p^{2}}{(1-u_{n})^2} \left(\frac{w_{n-1}}{w_{n}} \right) \left(1 - \frac{w_{n-1}}{w_{n}} \right) \left( \mathbf{P}_{n-1} - \mathbf{P}_{n} \right)
    \end{align*}
    \vspace{3pt}
    \end{minipage}
    }

    The curvature of a clamped NURBS curve at its endpoints is given by:
    \begin{align*}
        \kappa(u=0) &=  \left( \frac{p-1}{p}  \right)  \left( \frac{u_{p+1}}{u_{p+2}}  \right) \left(\frac{w_{0}\,w_{2}}{w_{1}^2} \right)  \,
        \frac{\norm{ \left(\mathbf{P}_{2}-\mathbf{P}_{0}\right) \times \left(\mathbf{P}_{1}-\mathbf{P}_{0}\right)} }{\norm{\mathbf{P}_{1}-\mathbf{P}_{0}}^3} \\
        \kappa(u=1) &=  \left( \frac{p-1}{p}  \right)  \left( \frac{1-u_{n}}{1-u_{n-1}}  \right) \left(\frac{w_{n}\,w_{n-2}}{w_{n-1}^2} \right)  \,
        \frac{\norm{ \left(\mathbf{P}_{n-2}-\mathbf{P}_{n}\right) \times \left(\mathbf{P}_{n-1}-\mathbf{P}_{n}\right)} }{\norm{\mathbf{P}_{n-1}-\mathbf{P}_{n}}^3}
    \end{align*}


    
\end{enumerate}



%% --------------------------------------------------------------------------------------------- %%
%% --------------------------------------------------------------------------------------------- %%
%% --------------------------------------------------------------------------------------------- %%



\clearpage

\section{NURBS surfaces}

%% --------------------------------------------------------------------------------------------- %%
%% --------------------------------------------------------------------------------------------- %%
%% --------------------------------------------------------------------------------------------- %%
\subsection{Definition}

A Non-Uniform Rational Basis Spline (NURBS) surface is a parametric surface defined by:
\begin{align}
\label{eq:def_nurbs_surface}
\mathbf{S}(u,v)=\sum_{i=0}^{n} \sum_{j=0}^{m} R_{i,j}^{p, q}(u, v) \, \mathbf{P}_{i,j} \quad \quad 0 \leq (u, \, v) \leq 1
\end{align}

where $p$ and $q$ are the orders of the surface in the $u$- and $v$-directions, the coefficients $\mathbf{P}_{i, j}$ are a bidirectional net of control points and $R_{i,j}^{p, q}(u, v)$ are the rational basis functions given by:
\begin{align}
    R_{i,j}^{p, q}(u, v) = \frac{N_{i,p}(u)N_{j,q}(v) \, w_{i,j}}{\sum\limits_{k=0}^{n} \sum\limits_{l=0}^{m} N_{k,p}(u) N_{l,p}(v) \, w_{k, l}}
\end{align}

in which $w_{i, j}$ are the weights of the control points and $N_{i,p}(u)N_{j,q}(v)$ are the product of univariate B-Spline basis functions defined on the non-decreasing knot vectors $U$ and $V$:
\begin{align}
    U = [u_0,\, \dots,\, u_{r}] \in \mathds{R}^{r+1} \quad &\text{with} \quad r = n + p + 1 \\
    V = [v_0,\, \dots,\, v_{s}] \in \mathds{R}^{s+1} \quad &\text{with} \quad s = m + q + 1
\end{align}

The $u$-direction basis functions are given by the recursive relation:
\begin{align}
    N_{i,0}(u) =  \begin{cases}
    1& \mathrm{if} \quad u_{i}\leq u<u_{i+1}\\
    0& \mathrm{otherwise}
    \end{cases}
\end{align}
\begin{align}
\label{eq:def_bspline_recursive_2}
 N_{i,p}(u)=\frac {u-u_{i}}{u_{i+p}-u_{i}} N_{i,p-1}(u) + \frac {u_{i+p+1}-u}{u_{i+p+1}-u_{i+1}} N_{i+1,p-1}(u)
\end{align}

whereas the $v$-direction basis functions are defined in an analogous way replacing the variable $u$ by $v$ and the indices $i$ and $p$ by $j$ and $q$, respectively.




%% --------------------------------------------------------------------------------------------- %%
%% --------------------------------------------------------------------------------------------- %%
%% --------------------------------------------------------------------------------------------- %%
\subsection{Mathematical properties of bivariate rational basis functions}
Here is a list of some important properties of the bivariate rational basis functions:

\begin{enumerate}
    
    \item The relation $r = n + p + 1$ holds, where $n+1$ is the number of basis functions in the $u$-direction and $r+1$ is the number of elements of the knot vector $U$.

    The relation $s = m + q + 1$ holds, where $m+1$ is the number of basis functions in the $v$-direction and $s+1$ is the number of elements of the knot vector $V$.
    
    \item The numerator and denominator of $R_{i,j}^{p, q}(u, v)$ are, at most, polynomials of degree $p \times q$.
    
    \item Local support:
    \begin{enumerate}
        \item $R_{i,j}^{p, q}(u, v)=0$ if $(u,v)$ is outside the rectangle $[u_{i}, u_{i+p+1}) \times [v_{j}, v_{j+q+1})$.
        \item In any given knot rectangle $[u_{i_{0}}, u_{i_{0}+1}) \times [v_{j_{0}}, v_{j_{0}+1})$, at most $(p+1)(q+1)$ basis functions are nonzero, namely the $R_{i,j}^{p, q}(u, v)$ with $i_{0}-p \leq i \leq i_{0}$ and $j_{0}-q \leq j \leq j_{0}$.
    \end{enumerate}
    
    \item Non-negativity:
        \begin{equation*}
            R_{i,j}^{p, q}(u, v) \geq 0 \quad \text{for all $i$, $j$, $p$, $q$, and $(u,v) \in [0, 1] \times [0, 1]$}
        \end{equation*}

    \item Partition of unity:
        \begin{equation*}
            \sum_{i=0}^n  \sum_{j=0}^{m} R_{i,j}^{p, q}(u, v) = 1 \text{ for all } (u,v) \in [0, 1] \times [0, 1]
        \end{equation*}
        
        In addition, for an arbitrary knot rectangle, $[u_{i_{0}}, u_{i_{0}+1}) \times [v_{j_{0}}, v_{j_{0}+1})$ we have that:
            \begin{equation*}
                \sum_{i=i_{0}-p}^{i_{0}}\sum_{j=j_{0}-q}^{j_{0}} R_{i,j}^{p, q}(u, v) = 1 \text{ for all } (u,v) \in [u_{i_{0}}, u_{i_{0}+1}) \times [v_{j_{0}}, v_{j_{0}+1})
            \end{equation*}
        This means that the sum of the non-zero basis-functions of any knot rectangle is unity.
        
    \item Continuity and differentiability:
    \begin{enumerate}
        \item The basis functions are infinitely differentiable in the interior of the knot rectangles formed by the $U$ and $V$ vectors.
        \item The basis functions are $p-k$ ($q-k$) continuously differentiable in the $u$-direction ($v$-direction) at a $u$-knot ($v$-knot) with multiplicity $k$.
    \end{enumerate}

    \item Extrema:
    
    $R_{i,j}^{p, q}(u, v)$ attains exactly one maximum value in the region $(u,v) \in [0, 1] \times [0, 1]$.
    
    \item There is no compact and useful expression for the first partial derivatives of the basis functions.
    
    \item There is no compact and useful expression for the $(k,l)$-th order partial derivatives of the basis functions.
    
    \item If all the control point weights are equal, then rational basis functions reduce to B-Spline basis functions, that is, $R_{i,j}^{p, q}(u, v)=N_{i,p}(u)N_{j,q}(v)$.
    
    \item If all the control point weights are equal, the knot vectors are clamped, and $(p,q)=(n,m)$, then rational basis function reduce to Bernstein polynomials, that is, $R_{i,j}^{p, q}(u, v)=B_{i,n}(u)B_{j,m}(v)$.
    
\end{enumerate}



%% --------------------------------------------------------------------------------------------- %%
%% --------------------------------------------------------------------------------------------- %%
%% --------------------------------------------------------------------------------------------- %%
\subsection{Mathematical properties of NURBS surface}
Here is a list of some important properties of NURBS surfaces:

\begin{enumerate}

    \item $\mathbf{S}(u, v)$ is a piecewise surface and its components are ratios of bivariate polynomials of, at most, degree $p\times q$.
    
    \item In the $u$-direction: the degree $p$, the number of control points $n+1$ and the number of knots $r+1$ are related according to $r = n + p + 1$
    
    In the $v$-direction: the degree $q$, the number of control points $m+1$ and the number of knots $s+1$ are related according to $s = m + q + 1$
    
    \item NURBS surfaces contain B-Spline and rational/non-rational \Bezier surfaces as special cases:
    \begin{enumerate}
        \item If all the control point weights are equal then the NURBS surface is reduced to a B-Spline surface.
        \item If $(p,q) = (n,m)$ and the knot vectors are clamped then the NURBS surface is reduced to a rational \Bezier surface.
        \item If all the control point weights are equal, $(p,q) = (n,m)$, and the knot vectors are clamped then the NURBS surface is reduced to a polynomial \Bezier surface.
    \end{enumerate}
    
    \item Affine invariance:
    
    NURBS surfaces are invariant under affine transformations such as rotations, displacements and scalings. This means that one can apply an affine transformation to the NURBS surface by applying it to its set of control points.
    
    Given an affine transformation $\phi(\mathbf{v}) = \mathrm{A} \mathbf{v} + \mathbf{b}$ we have that:
    \begin{align*}
        \phi \big(\mathbf{S}(u,v)\big) &= \phi \big(\sum_{i=0}^{n}\sum_{j=0}^{m} R_{i,j}^{p, q}(u, v)\, \mathbf{P}_{i,j}\big) = \\
        &= \mathrm{A} \sum_{i=0}^{n}\sum_{j=0}^{m} R_{i,j}^{p, q}(u, v)\, \mathbf{P}_{i,j} + \mathbf{b} = \\
        &= \mathrm{A} \sum_{i=0}^{n}\sum_{j=0}^{m} R_{i,j}^{p, q}(u, v)\, \mathbf{P}_{i,j} + \mathbf{b} \, \sum_{i=0}^{n}\sum_{j=0}^{m} R_{i,j}^{p, q}(u, v) \qquad \text{(by partition of unity)}\\
        &= \sum_{i=0}^{n}\sum_{j=0}^{m} R_{i,j}^{p, q}(u, v)\,(\mathrm{A} \mathbf{P}_{i,j} + \mathbf{b}) \\
        &= \sum_{i=0}^{n}\sum_{j=0}^{m} R_{i,j}^{p, q}(u, v)\,\phi \big(\mathbf{P}_{i,j} \big)
    \end{align*}
    
    \item Convex hull property:
    
    All the points of a NURBS surface are contained in the \textit{convex hull} of its control points.
    The convex hull of a set of points $P = \{\mathbf{P}_{0},\, \dots,\, \mathbf{P}_{N}\}$ is denoted by $\mathcal{CH}(P)$ and it is the set of all the convex combinations of points:
    \begin{equation*}
     \mathcal{CH}(P) = \Big\{ \mathbf{C} =  \sum_{k=0}^{N} a_{k} \, \mathbf{P}_{k} \text{ such that } \sum_{k=0}^{N} a_{k} = 1 \text{ and } a_{k}\geq 0 \text{ for } k = 0,\, 1,\, \dots,\, N\Big\}
    \end{equation*}
    The convex hull property follows from the non-negativity and the partition-of-unity properties of the basis functions.
    % The convex hull of a set of points is the minimal convex polytope that contains all the points
    
    \item Strong convex hull property:
    
    If $(u,v) \in [u_{i_{0}},\, u_{i_{0}+1}) \times [v_{j_{0}},\, v_{j_{0}+1})$, then $\mathbf{S}(u, v)$ is contained in the convex hull of the control points $\boldsymbol{P}_{i,j}$ with $i_{0}-p \leq i \leq i_{0}$ and $j_{0}-q \leq j \leq j_{0}$.

    \item Local modification scheme:
    
    Modifying the control point $\mathbf{P}_{i,j}$ or weight $w_{i,j}$ affects $\mathbf{S}(u, v)$ only in the rectangle $[u_{i}, u_{i+p+1}) \times [v_{j}, v_{j+q+1})$. This property follows from the fact that $ R_{i,j}^{p, q}(u, v)=0$ if $(u,v)$ is outside the rectangle $[u_{i}, u_{i+p+1}) \times [v_{j}, v_{j+q+1})$ and it implies that the shape of a NURBS can be modified locally without changing its shape globally.
    
    \item If triangulated, the net of control points represents a piecewise planar approximation to the NURBS surface.

    \item No known variation diminishing property.
    
    \item Continuity and differentiability:
    \begin{enumerate}
        \item NURBS surfaces are infinitely differentiable in the interior of the knot rectangles formed by the $U$ and $V$ vectors.
        \item NURBS surfaces are $p-k$ ($q-k$) continuously differentiable in the $u$-direction ($v$-direction) at a $u$-knot ($v$-knot) with multiplicity $k$.
    \end{enumerate}
    
    
    \item Homogeneous coordinates:
    
    $N$-dimensional rational functions with a common denominator can be represented as polynomial functions in $(N+1)$-dimensional space using \textit{homogeneous coordinates}.
    
    Consider a three-dimensional point $\mathbf{P}=(x,\, y,\, z)$. $\mathbf{P}$ can be written in four-dimensional space as $\mathbf{P}^w=(wx,\, wy,\, wz,\, w) = (X,\, Y,\, Z,\, W)$, where $w \neq 0$.
    The point $\mathbf{P}$ can be obtained from $\mathbf{P}^w$ using the mapping $\mathcal{H}$ given by:
    \begin{align*}
        \mathbf{P} = \mathcal{H}\{\mathbf{P}^w\} =  \mathcal{H}\{(X,\, Y,\, Z,\, W)\} = \left( \frac{X}{W},\, \frac{Y}{W},\, \frac{Z}{W}\right)
    \end{align*}
    
    This theory can be used to represent a $N$-dimensional NURBS surface as a $(N+1)$-dimensional B-Spline surface and then retrieve the coordinates of the NURBS surface using the $\mathcal{H}$ mapping.
    
    Indeed, consider a NURBS surface $\mathbf{S}(u, v)$ with control points $\mathbf{P}_{i, j}$ and weights $w_{i, j}$. It is possible to construct a set of weighted control points $\mathbf{P}^{w}_{i,j} = (w_{i,j}x_{i,j},\, w_{i,j}y_{i,j},\, w_{i,j}z_{i,j},\, w_{i,j})$ and define the corresponding non-rational B-Spline surface in four-dimensional space as:
    \begin{align*}
        \mathbf{S}^{w}(u,v)=\sum_{i=0}^{n} \sum_{j=0}^{m} R_{i,j}^{p,q}(u,v)\, \mathbf{P}^{w}_{i,j} \quad \quad 0 \leq (u,v) \leq 1
    \end{align*}
    
    In order to compute the coordinates of the original NURBS surface it is possible to use standard B-Spline algorithms to evaluate the coordinates of $\mathbf{S}^{w}(u,v)$ in homogeneous space and then apply the $\mathcal{H}$ mapping:
    \begin{align*}
        \mathbf{S}(u,v) &= \mathcal{H}\{\mathbf{S}^w\} \\
        &=  \mathcal{H} \bigg\{\sum_{i=0}^{n} \sum_{j=0}^{m} N_{i,p}(u)N_{j,q}(u,v)\, \mathbf{P}^{w}_{i,j} \bigg\}  \\
        &= \frac{\sum\limits_{i=0}^{n} \sum\limits_{j=0}^{m} N_{i,p}(u)N_{j,q}(u)\, w_{i,j} \,\mathbf{P}_{i,j}}{\sum\limits_{i=0}^{n} \sum\limits_{j=0}^{m} N_{i,p}(u)N_{j,q}(v)\, w_{i,j}} \\
        &= \sum_{i=0}^{n}\sum_{j=0}^{m} R_{i,j}^{p,q}(u,v) \, \mathbf{P}_{i,j}
    \end{align*}
    
    This property is important to evaluate the derivatives of NURBS surfaces.
    
    \item First and higher order derivatives:
    
    The derivatives of a NURBS surface $\mathbf{S}(u,v)$ can be expressed in terms of the derivatives of the corresponding B-Spline in homogeneous space $\mathbf{S}^{w}(u,v)$. Indeed, let:
    \begin{align*}
        \mathbf{S}(u,v) = \frac{\sum\limits_{i=0}^{n}\sum\limits_{j=0}^{m} N_{i,p}(u)N_{j,q}(v)\, w_{i,j} \,\mathbf{P}_{i,j}}{\sum\limits_{i=0}^{n} \sum\limits_{j=0}^{m} N_{i,p}(u)N_{j,q}(u)\, w_{i,j}} = \frac{\mathbf{A}(u,v)}{w(u,v)}
    \end{align*}
    
    where $\mathbf{A}(u,v)$ is the vector function whose coordinates are the three first elements of $\mathbf{S}^{w}(u,v)$ and $w(u,v)$ is the scalar function with the fourth element of $\mathbf{S}^{w}(u,v)$.
    
    To compute the first partial derivatives, first rearrange the previous expression to avoid the denominator, then differentiate both sides of the equation and, finally, solve for $\frac{\partial\mathbf{S}}{\partial \alpha}$, where $\alpha$ is either $u$ or $v$:
    \begin{align*}
        & w \, \mathbf{S} = \mathbf{A} \\
        & \big[ w \, \mathbf{S} \big]_{\alpha}  = \mathbf{A}_{\alpha} \\
        & w_{\alpha} \, \mathbf{S} + w \, \mathbf{S}_{\alpha} = \mathbf{A}_{\alpha} \\
        & \frac{\partial \mathbf{S}}{\partial \alpha} = \mathbf{S}_{\alpha} = \frac{1}{w} \left( \mathbf{A}_{\alpha} - w_{\alpha} \, \mathbf{S} \right)
    \end{align*}
    
    To compute the $(k,l)$-th order partial derivatives, differentiate $(k,l)$-times using Leibniz' rule for product differentiation and then solve for $\mathbf{S}^{(k)}(u,v)$: 
    
    \resizebox{\linewidth}{!}{
    \begin{minipage}{0.99\linewidth}
    \begin{align*}
        & w \, \mathbf{S} = \mathbf{A} \\
        & \big[ w \, \mathbf{S} \big]^{(k,l)} = \mathbf{A}^{(k,l)} \\
        & \sum_{i=0}^{k}\sum_{j=0}^{l} \binom{k}{i}\binom{l}{j} w^{(i,j)} \, \mathbf{S}^{(k-i, l-j)} = \mathbf{A}^{(k, l)} \\
        & \mathbf{S}^{(k,l)} = \frac{1}{w} \left[ \mathbf{A}^{(k,l)} - \sum\limits_{i=1}^{k} \binom{k}{i} w^{(i,0)} \, \mathbf{S}^{(k-i,l)} - \sum\limits_{j=1}^{l} \binom{l}{j} w^{(0,j)} \, \mathbf{S}^{(k,l-j)} - \sum\limits_{i=1}^{k}\sum\limits_{j=1}^{l} \binom{k}{i}\binom{l}{j} w^{(i,j)} \, \mathbf{S}^{(k-i, l-j)} \right]
    \end{align*}
    \vspace{3pt}
    \end{minipage}
    }
    
    The derivatives of $\mathbf{A}(u,v)$ and $w(u,v)$ can be easily computed by differentiating $\mathbf{S}^{w}(u,v)$:
        \begin{align*}
             \left[\mathbf{S}^{w}\right]^{(k,l)}(u,v)  = \sum_{i=0}^{n}\sum_{j=0}^{m} N^{(k)}_{i,p}(u)N^{(l)}_{j,q}(u)\, \mathbf{P}^{w}_{i,j}
        \end{align*}
        
    \item Corner point interpolation:
    
    The corners of a clamped NURBS surface coincide with the corner points of its control net:
    \begin{align*}
        \mathbf{S}(u=0,\, v=0) &= \mathbf{P}_{0, 0} \\
        \mathbf{S}(u=1,\, v=0) &= \mathbf{P}_{n, 0} \\
        \mathbf{S}(u=0,\, v=1) &= \mathbf{P}_{0, m} \\
        \mathbf{S}(u=1,\, v=1) &= \mathbf{P}_{n, m} \\
    \end{align*}
    

    
\end{enumerate}


%% --------------------------------------------------------------------------------------------- %%
%% --------------------------------------------------------------------------------------------- %%
%% --------------------------------------------------------------------------------------------- %%



 \printbibliography



\end{document}